\section{Experimental setup}
\label{sec:setup}

\subsection{Machines}

All three machines are personally owned consumer laptops on one domestic Wi-Fi
network. Nothing was rented, and no datacentre hardware was involved.

\begin{table}[htbp]
\centering\small
\caption{Hardware. Video memory figures are as reported by the machine itself
through the sources in Table~\ref{tab:sources}.}
\label{tab:hardware}
\begin{tabular}{@{}lllrr@{}}
\toprule
\textbf{Role} & \textbf{Machine} & \textbf{Accelerator} & \textbf{GPU memory} & \textbf{System RAM} \\
\midrule
Controller      & ASUS ROG, Linux   & RTX 3050 Laptop, CUDA & \SI{4096}{MiB} & \SI{15234}{MiB} \\
Remote worker   & ASUS TUF, Windows & RTX 4050 Laptop, CUDA & \SI{6141}{MiB} & \SI{16003}{MiB} \\
Remote worker   & MacBook Air, macOS & Apple M4, Metal      & unified        & \SI{16384}{MiB} \\
\bottomrule
\end{tabular}
\end{table}

The Mac row carries no separate GPU memory figure because Apple Silicon has no
separate pool to report. Its \SI{16384}{MiB} is one unified allocation shared
between CPU and GPU.

\subsection{Software}

Both NVIDIA machines run \texttt{llama.cpp} release \code{b11909}, built with
CUDA and RPC support, at RPC protocol 8.0.1. The Mac runs the official
\code{b11909} macOS arm64 build. Keeping one release across all three machines
removed protocol mismatch as a variable. Every remote
\code{ggml-rpc-server} runs with \code{-c}, which caches received tensors on the
worker's own disk so that only the first load transfers weights over Wi-Fi.

\subsection{Models}

\begin{itemize}[leftmargin=1.4em,itemsep=2pt]
  \item \textbf{Llama 3.2 1B Instruct}, Q4\_K\_M --- local fallback profile.
  \item \textbf{Llama 3.2 3B Instruct}, Q4\_K\_M, \SI{2.0}{GB} --- used for every
        split and recovery experiment. It fits in the controller's GPU alone,
        which is what allows the controller to take over after a worker is lost.
  \item \textbf{Llama 3.1 8B Instruct}, Q4\_K\_M, \SI{4.9}{GB} --- capacity
        experiment. Verified byte-for-byte against the published artefact. It
        does not fit in the controller's \SI{4096}{MiB} GPU alone.
\end{itemize}

No weights are committed to the repository. Context length was 2{,}048 tokens
throughout.

\subsection{Network}

A \SI{5}{GHz} Wi-Fi network. Link conditions were recorded at measurement time
because they are the dominant variable in every distributed figure here. During
the benchmark runs of Section~\ref{sec:perf} the controller reported
\SI{-56}{dBm} with driver-reported rates of \SI{780}{Mbit/s} receive and
\SI{292.5}{Mbit/s} transmit. During the earliest model transfers, several days
before, the same link reported \SI{-66}{dBm} to \SI{-68}{dBm} and
\SI{26}{Mbit/s} to \SI{29}{Mbit/s} transmit. These are physical-layer rates
negotiated by the driver, not achieved application throughput, and the gap
between the two sessions explains why first-load times differ so widely in
Section~\ref{sec:load}.

\subsection{Method}

Performance figures come from five sequential requests against an
already-loaded model, using one fixed prompt asking for a 100-word explanation
of distributed computing, capped at 160 output tokens. Normal sampling was left
enabled, so output lengths vary slightly between runs; we therefore report
medians and full ranges rather than a single best run. End-to-end latency is
\code{curl}'s \code{time\_total}. Decode throughput is taken from
\texttt{llama.cpp}'s own per-request timing line, using the \code{eval time}
figure and excluding \code{prompt eval time}; it is never derived from response
character counts.

Three configurations were measured:

\begin{enumerate}[leftmargin=1.4em,itemsep=2pt]
  \item \textbf{Local baseline} --- \texttt{llama.cpp} on the controller alone,
        no RPC.
  \item \textbf{Native split} --- \texttt{llama.cpp} directly, with the Windows
        RPC endpoint and the same tensor weights HyperPod had computed.
  \item \textbf{HyperPod split} --- the same split reached through the Go
        controller.
\end{enumerate}

Comparing the second and third isolates the control plane, since the inference
path is identical. Failure experiments were performed by stopping a
\code{ggml-rpc-server} process by hand while a request was in flight; no
failures were simulated in software for the results reported here.
